\section*{अध्याय \ref{ch:g7:areas} --- क्षेत्रफल और आयतन}

\begin{solution}{exo:g7:areas:1}
$8 \times 5 = 40$ cm$^2$; \qquad $6.5 \times 4 = 26$ m$^2$.
\end{solution}

\begin{solution}{exo:g7:areas:2}
$A = 10 \times 4 = 40$ cm$^2$। उत्तर $10 \times 6 = 60$ नहीं है,
क्योंकि $6$ cm वाली भुजा तिरछी है: सूत्र \emph{ऊँचाई} (\omterm{def:g6:lines:perp}{लंब} दूरी)
माँगता है, जो $4$ cm है।
\end{solution}

\begin{solution}{exo:g7:areas:3}
$\frac{12 \times 7}{2} = 42$ cm$^2$; \quad
$\frac{9 \times 6}{2} = 27$; \quad
$\frac{5.5 \times 2}{2} = 5.5$ m$^2$.
\end{solution}

\begin{solution}{exo:g7:areas:4}
चक्र: $\pi \times 5^2 = 25\pi \approx 79$ cm$^2$। $10$ \omterm{def:g3:shapes:shapes}{त्रिज्या}
वाला \omterm{def:g3:division:half}{आधा} चक्र: $\frac{\pi \times 10^2}{2} = 50\pi \approx 157$ cm$^2$।
\end{solution}

\begin{solution}{exo:g7:areas:5}
\omterm{def:g6:measure:perimeter}{परिमाप}: $2\pi \times 3 = 6\pi \approx 18.8$ cm (एक लंबाई, cm में)।
\omterm{def:g6:measure:area}{क्षेत्रफल}: $\pi \times 3^2 = 9\pi \approx 28.3$ cm$^2$ (एक सतह,
cm$^2$ में)।
\end{solution}

\begin{solution}{exo:g7:areas:6}
\omterm{def:g7:areas:prism}{प्रिज़्म}: $24 \times 10 = 240$ cm$^3$। \omterm{def:g7:areas:prism}{बेलन}:
$\pi \times 3^2 \times 8 = 72\pi \approx 226$ cm$^3$।
\end{solution}

\begin{solution}{exo:g7:areas:7}
$\frac{8 \times h}{2} = 28$, इसलिए $4h = 28$ और $h = 7$ cm।
\end{solution}

\begin{solution}{exo:g7:areas:8}
तीनों \omterm{def:g2:mult:def}{गुणनफल} इसलिए मिलते हैं कि हर एक \emph{वही} राशि निकालता है
--- त्रिभुज का \omterm{def:g6:measure:area}{क्षेत्रफल}। इससे रचनाओं की एक काम की जाँच मिलती है,
और किसी दूसरे आधार-ऊँचाई जोड़े से कोई ऊँचाई निकालने का तरीक़ा भी
(\cref{exo:g7:areas:12} में इस्तेमाल हुआ)।
\end{solution}

\begin{solution}{exo:g7:areas:9}
दोनों त्रिभुजों के \omterm{def:g6:measure:area}{क्षेत्रफल} $\frac{30 \times 20}{2} = 300$ m$^2$
और $\frac{18 \times 20}{2} = 180$ m$^2$ हैं: कुल $480$ m$^2$। जो आम
सूत्र सूझता है: $h$ दूरी पर \omterm{def:g6:lines:perp}{समांतर} भुजाओं $a$ और $b$ के लिए,
\[
A = \frac{(a + b) \times h}{2}
\]
--- यानी समलंब का सूत्र (यहाँ $\frac{(30+18)\times 20}{2} = 480$)।
\end{solution}

\begin{solution}{exo:g7:areas:10}
पानी का \omterm{def:g6:measure:volume}{आयतन}: $\pi \times 6^2 \times 15 = 540\pi \approx 1\,696$
cm$^3$। डिब्बे का आधार: $18 \times 12 = 216$ cm$^2$। बनने वाली
ऊँचाई: $\frac{540\pi}{216} = 2.5\pi \approx 7.9$ cm।
\end{solution}

\begin{solution}{exo:g7:areas:11}
\omterm{def:g3:shapes:shapes}{त्रिज्याएँ} $15$ और $20$ cm। \omterm{def:g6:measure:area}{क्षेत्रफल}: $225\pi \approx 707$ cm$^2$
और $400\pi \approx 1257$ cm$^2$। $100$ cm$^2$ का दाम:
$\frac{9}{7.07} \approx 1.27$ यूरो और
$\frac{14}{12.57} \approx 1.11$ यूरो। बड़ा पिज़्ज़ा ज़्यादा फ़ायदे
का सौदा है --- \omterm{def:g6:measure:area}{क्षेत्रफल} \omterm{def:g4:geometry:circle}{व्यास} के \emph{वर्ग} के साथ बढ़ता है।
\end{solution}

\begin{solution}{exo:g7:areas:12}
समकोण-भुजाओं से: $A = \frac{6 \times 8}{2} = 24$ cm$^2$। कर्ण को
आधार मानकर: $24 = \frac{10 \times h}{2} = 5h$, इसलिए
$h = 4.8$ cm।
\end{solution}

\begin{solution}{pb:g7:areas:1}
\textbf{1.} वर्ग: $8 \times 8 = 64$। आयत: $5 \times 13 = 65$।
वही चारों टुकड़े एक आकृति में $64$ वर्ग इकाई ढकते दिखते हैं और
दूसरी में $65$: क़िस्सा यह है कि एक इकाई \omterm{def:g6:measure:area}{क्षेत्रफल} \omterm{def:g1:counting:zero}{शून्य} से पैदा हो
गया।

\textbf{2.} प्रत्येक त्रिभुज: $\frac{8 \times 3}{2} = 12$।
प्रत्येक समलंब: $\frac{(5 + 3) \times 5}{2} = 20$।

\textbf{3.} $12 + 12 + 20 + 20 = 64$: टुकड़ों का कुल $64$ है,
इसलिए वे वर्ग को ठीक-ठीक भर सकते हैं --- और $65$ इकाई वाले आयत को
भर \emph{नहीं} सकते। आयत वाली तस्वीर में कहीं न कहीं एक छेद होना ही
चाहिए।

\textbf{4.} एक जैसी ढलान का मतलब होता कि $8$ में $3$, $5$ में $2$
के \omterm{def:g6:propdata:def}{समानुपाती} हो; तिरछे \omterm{def:g2:mult:def}{गुणनफल}: $3 \times 5 = 15$ और
$2 \times 8 = 16$। बराबर नहीं: कर्ण और तिरछी किनार एक सीध में
\emph{नहीं} बैठते। आयत का ``विकर्ण'' एक झूठ है --- दो ज़रा-सी
अलग ढलानें, जो एक बहुत चौड़े कोण पर मिलती हैं।

\textbf{5.} उस झूठे विकर्ण के साथ चारों टुकड़े एक लंबी, बेहद पतली
ख़ाली जगह छोड़ देते हैं --- बहुत चपटे \omterm{def:g7:central:parallelogram}{समांतर चतुर्भुज} की शक्ल की एक
झिरी --- जिसका \omterm{def:g6:measure:area}{क्षेत्रफल} ठीक वही ग़ायब $65 - 64 = 1$ वर्ग इकाई है।
क़रीब $13$ इकाई लंबे विकर्ण पर फैलने पर उसकी औसत चौड़ाई क़रीब
$1 \div 13 \approx 0.08$ इकाई निकलती है: असली चित्र में उस पेंसिल की
लकीर से भी पतली जो उसे छिपाए रखती है।

\textbf{6.} प्रत्येक त्रिभुज का आधार $6$ cm और ऊँचाई $4$ cm है ---
यानी दोनों \omterm{def:g6:lines:perp}{समांतर} रेखाओं के बीच की दूरी, शिखर चाहे जहाँ बैठे ---
इसलिए हर एक का \omterm{def:g6:measure:area}{क्षेत्रफल} $\frac{6 \times 4}{2} = 12$ cm$^2$ है
(\cref{thm:g7:areas:triangle})। शिखर को \omterm{def:g6:lines:perp}{समांतर} रेखा पर सरकाने से
आकार बदलता है, आधार या ऊँचाई कभी नहीं: \omterm{def:g6:measure:area}{क्षेत्रफल} सदा बराबर।

\textbf{7.} विकर्ण समलंब को आधार $B$ वाले एक त्रिभुज और आधार $b$
वाले दूसरे त्रिभुज में काट देता है, और दोनों की ऊँचाई $h$ है (यानी
\omterm{def:g6:lines:perp}{समांतर} भुजाओं के बीच की दूरी):
\[
A = \frac{B \times h}{2} + \frac{b \times h}{2}
= \frac{(B + b) \times h}{2} .
\]

\textbf{8.} $A = \frac{(9 + 5) \times 4}{2} = \frac{56}{2} =
28$ cm$^2$।

\textbf{9.} \omterm{def:g6:lines:perp}{समांतर} भुजाओं का औसत $\frac{B + b}{2}$ है, इसलिए
(औसत) $\times h = \frac{(B+b) \times h}{2}$: वही सूत्र, बस अलग
तरह से गुणनखंडों में लिखा हुआ। जाँच: औसत $= \frac{9+5}{2} = 7$, और
$7 \times 4 = 28$ cm$^2$।

\textbf{10.} तिरछी ढेरी में ठीक वही पत्ते हैं --- वही परतें, हर एक
उतने ही \omterm{def:g6:measure:area}{क्षेत्रफल} की --- इसलिए \omterm{def:g6:measure:volume}{आयतन} भी वही है। $V = B \times h$ में
ऊँचाई आधार के \emph{\omterm{def:g6:lines:perp}{लंब}} नापी जानी चाहिए (ढेरी की खड़ी ऊँचाई),
तिरछाई के साथ नहीं: ठीक वैसे ही जैसे \omterm{def:g7:central:parallelogram}{समांतर चतुर्भुज} का \omterm{def:g6:measure:area}{क्षेत्रफल}
\omterm{def:g6:lines:perp}{लंब} ऊँचाई माँगता है, तिरछी भुजा नहीं (\cref{exo:g7:areas:2}) ---
बस एक विमा ऊपर।

\textbf{11.} छल्ला $= \pi \times 5^2 - \pi \times 3^2 = 25\pi -
9\pi = 16\pi \approx 50$ m$^2$ (\cref{thm:g7:areas:disk})।

\textbf{12.} $16\pi$ \omterm{def:g6:measure:area}{क्षेत्रफल} वाले चक्र की \omterm{def:g3:shapes:shapes}{त्रिज्या} $4$ है,
क्योंकि $\pi \times 4^2 = 16\pi$। \omterm{def:g3:shapes:shapes}{त्रिज्याएँ} $3$, $4$, $5$: गणित की
सबसे मशहूर तिकड़ी, जो \cref{ch:g8:pythagoras} में मुख्य भूमिका
निभाती है।

\textbf{13.} $40$ cm वाला एक पिज़्ज़ा: \omterm{def:g3:shapes:shapes}{त्रिज्या} $20$, \omterm{def:g6:measure:area}{क्षेत्रफल}
$\pi \times 400 \approx 1\,257$ cm$^2$। $28$ cm वाले दो पिज़्ज़े:
\omterm{def:g3:shapes:shapes}{त्रिज्या} $14$, दोनों मिलाकर
$2 \times \pi \times 196 = 392\pi \approx 1\,232$ cm$^2$। अकेला बड़ा
पिज़्ज़ा दोनों मँझोलों से \emph{ज़्यादा} पिज़्ज़ा है।

\textbf{14.} \omterm{def:g3:shapes:shapes}{त्रिज्या} को $1.4$ से बढ़ाने पर \omterm{def:g6:measure:area}{क्षेत्रफल}
$1.4^2 = 1.96$ गुना हो जाता है --- लगभग दुगुना। ठीक-ठीक गुणक वह
संख्या है जिसका वर्ग $2$ है, क़रीब $1.414$: यही
\cref{ex:g8:pythagoras:diagonal} वाली विकर्ण-संख्या है। (मोटा
नियम: डेढ़ गुनी चौड़ाई माने लगभग दुगुना पिज़्ज़ा।)

\textbf{15.} $13 \times 13 = 169$, पर $8 \times 21 = 168$: इस बार
एक वर्ग इकाई \emph{ग़ायब} हो जाती है --- टुकड़े ख़ाली जगह छोड़ने के
बजाय एक पतली झिरी में एक-दूसरे पर चढ़ जाते हैं।
हेमचंद्र--फ़िबोनाच्ची श्रेणी के लगातार पद $3, 5, 8, 13, 21$ लेने पर
\omterm{def:g2:mult:def}{गुणनफल} बारी-बारी से वर्ग से एक ज़्यादा और एक कम निकलते हैं
($5 \times 13 = 65 = 8^2 + 1$, फिर $8 \times 21 = 168 = 13^2 - 1$):
जादूगर बारी-बारी से एक इकाई ``कमाता'' और ``गँवाता'' है। \omterm{def:g6:measure:area}{क्षेत्रफल}
कभी झूठ नहीं बोलते: जिन टुकड़ों का कुल $64$ है वे $64$ ही ढकेंगे,
चाहे जहाँ सरकाओ --- हर ग़ायब या ज़्यादा इकाई किसी झिरी में छिपी
बैठी है, इस इंतज़ार में कि कोई ढलान जाँच ले।

\end{solution}
